\subsubsection{Anomaly Detection for Minimizing Epistemic Uncertainty}\label{sec:uncertainty}

For effective navigation in off-road environments, the traversability of terrain can be determined by interactions between terrain and a vehicle~\cite{wellhausen2019should,Z_acoustic}. However, data on such interactions is restricted to traversed regions only, and data collection in non-traversable regions is not feasible. This nature produces overconfident predictions, resulting in navigational failure or even catastrophic outcomes~\cite{Wellhausen_2020}.

This issue can be mitigated by employing methods for detecting anomalies~\cite{kerner2019novelty,Wellhausen_2020,pmlr-v155-ji21a}. Regions that a vehicle interacted with are designated as positive, on the assumption that terrains with similar geometric properties to those regions are normal and thus less uncertain. Then, we train a binary classifier capable of identifying points that are dissimilar to positive points, indicating high epistemic uncertainty. These regions are designated for avoidance during navigation. The method is described in detail in our previous work~\cite{seo_scate_2023, seo2023learning}.


